\documentclass[10pt,a4paper]{article}
\usepackage[utf8]{inputenc}
\usepackage[T1]{fontenc}
\usepackage{fouriernc}
\usepackage[french]{babel}
\frenchbsetup{StandardLists=true}
\NoAutoSpaceBeforeFDP
\usepackage[left=1.5cm,right=1.5cm,top=1.5cm,bottom=1.5cm]{geometry}
\usepackage{multirow,tabularx,booktabs,array}
\usepackage[dvipsnames]{xcolor}
\usepackage{amsmath}
\usepackage{fancyhdr}
\usepackage{colortbl}
\setlength{\headheight}{14pt}

\pagestyle{fancy}
\renewcommand\headrulewidth{1pt}
\fancyhead[L]{Lycée Denis Diderot}
\fancyhead[R]{Classe : T~Bac Pro}
\fancyfoot[C]{}
\fancyfoot[R]{\today}
\fancyfoot[L]{Alexis RUIZ}

\setlength{\parindent}{0pt}
\setlength{\extrarowheight}{0.4mm}
\renewcommand{\arraystretch}{1.35}

% Couleurs
\definecolor{exoheader}{RGB}{30,90,160}
\definecolor{exolight}{RGB}{220,235,255}
\definecolor{totalrow}{RGB}{255,230,200}

% Macro en-tête de partie
\newcommand{\exotitle}[2]{%
  \noindent\colorbox{exoheader}{%
    \parbox{\dimexpr\linewidth-2\fboxsep\relax}{%
      \textcolor{white}{\textbf{\large #1 \hfill #2}}%
    }%
  }%
}

\begin{document}

\begin{center}
  {\LARGE\textbf{GRILLE DE CORRECTION}}\\[1mm]
  {\large Moteurs électriques \hspace{1cm} T~Bac Pro \hspace{1cm} \textbf{Total~: /20 points}}
\end{center}

\vspace{3mm}

% ===================== PARTIE A =====================
\exotitle{Partie A : Identification des moteurs}{/4 pts}

\vspace{1mm}

\begin{tabularx}{\linewidth}{|c|X|c|}
  \hline
  \rowcolor{exolight}
  \textbf{Q} & \textbf{Réponse attendue} & \textbf{Pts} \\
  \hline
  \multirow{4}{*}{1} & M1 identifié : \textbf{moteur à courant continu (MCC)} & 0,5 \\
  \cline{2-3}
  & Justification : plaque indique \textbf{24\,V DC} (tension continue) & 0,5 \\
  \cline{2-3}
  & M2 identifié : \textbf{moteur asynchrone triphasé} & 0,5 \\
  \cline{2-3}
  & Justification : plaque indique \textbf{AC} et fréquence \textbf{f = 50\,Hz} & 0,5 \\
  \hline
  \multirow{3}{*}{2} & Calcul du nombre de paires de pôles : $p = \dfrac{60 \times 50}{3000} = 1$ & 0,5 \\
  \cline{2-3}
  & Formule et calcul : $g = \dfrac{3000 - 2840}{3000}$ & 0,5 \\
  \cline{2-3}
  & \textbf{Résultat encadré :} $g \approx 5{,}3~\%$ & 1 \\
  \hline
  \rowcolor{totalrow}
  \multicolumn{2}{|r|}{\textbf{Sous-total Partie A}} & \textbf{/4} \\
  \hline
\end{tabularx}

\vspace{3mm}

% ===================== PARTIE B =====================
\exotitle{Partie B : Bilan de puissance et rendement}{/6 pts}

\vspace{1mm}

\begin{tabularx}{\linewidth}{|c|X|c|}
  \hline
  \rowcolor{exolight}
  \textbf{Q} & \textbf{Réponse attendue} & \textbf{Pts} \\
  \hline
  \multirow{2}{*}{3} & Formule : $P_{\rm élec} = U \times I$ & 0,5 \\
  \cline{2-3}
  & Application + résultat encadré : $P_{\rm élec} = 24 \times 3{,}8 = \mathbf{91{,}2~W}$ & 0,5 \\
  \hline
  \multirow{4}{*}{4} & Formule : $\omega = \dfrac{2\pi\,n}{60}$ & 0,5 \\
  \cline{2-3}
  & Application + résultat : $\omega \approx 148{,}7~\text{rad/s}$ & 0,5 \\
  \cline{2-3}
  & Formule : $P_{\rm méc} = C \times \omega$ & 0,5 \\
  \cline{2-3}
  & Application + résultat encadré : $P_{\rm méc} = 0{,}55 \times 148{,}7 \approx \mathbf{81{,}8~W}$ & 0,5 \\
  \hline
  \multirow{4}{*}{5} & Calcul : $P_{\rm pertes} = P_{\rm élec} - P_{\rm méc} = 91{,}2 - 81{,}8 \approx \mathbf{9{,}4~W}$ & 0,75 \\
  \cline{2-3}
  & Calcul : $\eta = P_{\rm méc} / P_{\rm élec} \approx \mathbf{89{,}7~\%}$ & 0,75 \\
  \cline{2-3}
  & Estimation plaque : $\eta_{\rm plaque} = 90 / (24 \times 4{,}2) \approx 89{,}3~\%$ & 0,75 \\
  \cline{2-3}
  & Commentaire cohérent (valeurs proches, moteur légèrement sous charge nominale) & 0,75 \\
  \hline
  \rowcolor{totalrow}
  \multicolumn{2}{|r|}{\textbf{Sous-total Partie B}} & \textbf{/6} \\
  \hline
\end{tabularx}

\vspace{3mm}

% ===================== PARTIE C =====================
\exotitle{Partie C : Influence des paramètres électriques sur la vitesse}{/10 pts}

\vspace{1mm}

\begin{tabularx}{\linewidth}{|c|X|c|}
  \hline
  \rowcolor{exolight}
  \textbf{Q} & \textbf{Réponse attendue} & \textbf{Pts} \\
  \hline
  \multirow{3}{*}{6} & 6 points correctement placés dans le repère $n = f(U)$ & 1,5 \\
  \cline{2-3}
  & Droite de tendance tracée & 0,5 \\
  \cline{2-3}
  & Description : relation \textbf{linéaire / proportionnelle} entre $U$ et $n$ & 1 \\
  \hline
  7 & Conséquence correcte : faire varier $U$ permet de régler la vitesse de la perceuse (vitesse de coupe adaptable) & 1 \\
  \hline
  \multirow{3}{*}{8} & 5 points correctement placés dans le repère $n = f(f_{\rm réseau})$ & 1,5 \\
  \cline{2-3}
  & Droite de tendance tracée & 0,5 \\
  \cline{2-3}
  & Description : relation \textbf{linéaire / proportionnelle} entre $f$ et $n$ & 1 \\
  \hline
  \multirow{3}{*}{9} & Principe du variateur : convertit la fréquence réseau en fréquence réglable & 1 \\
  \cline{2-3}
  & Lien $n \propto f$ correctement évoqué & 1 \\
  \cline{2-3}
  & Exemple d'application industrielle cité (convoyeur, pompe, ventilateur…) & 1 \\
  \hline
  \rowcolor{totalrow}
  \multicolumn{2}{|r|}{\textbf{Sous-total Partie C}} & \textbf{/10} \\
  \hline
\end{tabularx}

\vspace{4mm}

\colorbox{exoheader}{%
  \parbox{\dimexpr\linewidth-2\fboxsep\relax}{%
    \textcolor{white}{\textbf{\large TOTAL GÉNÉRAL \hfill /20 points}}%
  }%
}

\end{document}
